\section{Lean bridge interfaces}

This chapter records the interface graph of the Lean bridge.  The
graph is the formal skeleton of the paper: it shows exactly which
statements are compiled, which are conditional, and which remain
open.

\subsection{Conventions}

\begin{itemize}[leftmargin=*]
\item \textbf{compiled}: the theorem compiles in the pinned Lean
  environment without \texttt{sorry}.
\item \textbf{compiled (conditional)}: the theorem compiles, but one
  of its inputs is an open proposition such as
  \texttt{unifiedCoreClosed} or \texttt{failureWindowExistence}.
\item \textbf{pending}: the theorem contains \texttt{sorry} or has not
  yet been assembled.
\end{itemize}

\subsection{Interface map}

\begin{verbatim}
FinitePrefix.lean (compiled)
  -> UnifiedCoreBridge.lean (compiled)
  -> UnifiedCoreAudit.lean (one pending: unified_core_final_no_hge)
  -> unifiedCoreClosed (open input)
  -> decisiveWindowValuationBoundCorrected_of_unified_core (compiled*)
  -> failureWindowExistenceOfUnifiedCore (compiled*)
  -> CycleBridge.windowBoundToNoCycle_of_failureWindowExistence
       (compiled*, conditional)
  -> DwdbDiv.dwdbDivFinalAssembly (compiled*, conditional)
  -> FinalStatement.five_x_plus_one_diverges_at_7 (pending)
\end{verbatim}

The arrows are implications in the theorem statements, not simulation
steps.  Each conditional node compiles because it is written as a
function of its open input.

\subsection{Conditional interfaces}

\begin{center}
\small
\begin{tabular}{@{}p{0.44\textwidth}p{0.30\textwidth}p{0.20\textwidth}@{}}
\toprule
Interface & Consumes & Status \\
\midrule
\texttt{decisiveWindowValuation\allowbreak
BoundCorrected\_\allowbreak of\_\allowbreak unified\_\allowbreak core}
& \texttt{unifiedCoreClosed} & compiled* \\
\texttt{failureWindowExistenceOfUnifiedCore} &
  \texttt{unifiedCoreClosed} & compiled* \\
\texttt{windowBoundToNoCycle\_\allowbreak of\_\allowbreak
failureWindowExistence} &
  \texttt{failureWindowExistence} & compiled* \\
\texttt{dwdbDivFinalAssembly} &
  \texttt{windowBoundToNoCycle} & compiled* \\
\texttt{five\_\allowbreak x\_\allowbreak plus\_\allowbreak one\_\allowbreak
diverges\_\allowbreak at\_\allowbreak 7} &
  \texttt{dwdbDivFinalAssembly} & pending \\
\bottomrule
\end{tabular}
\end{center}

\subsection{What is already closed}

The compiled layer includes:
\begin{itemize}[leftmargin=*]
\item the finite-prefix expansion and first-big-step facts;
\item the candidate parameterisation and reset-head predecessor;
\item the segment equations and the $d=1,d=2,d=3,d\ge4$ exclusions;
\item the CRT equation, the terminal residue equivalence, and the
  size-condition lower bound;
\item the cycle-word extraction and QB-8 split;
\item the failure-window contradiction lemma;
\item the window-to-no-cycle bridge, conditional on the open inputs.
\end{itemize}

\subsection{What remains open}

Exactly two statements are open:
\begin{center}
\begin{tabular}{@{}ll@{}}
\toprule
Statement & Status \\
\midrule
\texttt{UnifiedCoreAudit.unified\_\allowbreak core\_\allowbreak
final\_\allowbreak no\_\allowbreak hge} & pending \\
\texttt{FinalStatement.five\_\allowbreak x\_\allowbreak plus\_\allowbreak
one\_\allowbreak diverges\_\allowbreak at\_\allowbreak 7} & pending \\
\bottomrule
\end{tabular}
\end{center}

The second statement is an assembly target: once the first statement
is closed, the compiled conditional interfaces produce the final
theorem.

\subsection{Axiom audit}

The audit of the bridge theorems returns only
\begin{verbatim}
[propext, Classical.choice, Quot.sound]
\end{verbatim}
for theorems that use choice, and
\begin{verbatim}
[propext, Quot.sound]
\end{verbatim}
for theorems that do not.  The audit runs \texttt{\#print axioms} on
the actual theorems and compares the result with the allowlist.  No
custom axiom is used.

\begin{center}
\small
\begin{tabular}{@{}p{0.62\textwidth}p{0.32\textwidth}@{}}
\toprule
Theorem & Axioms \\
\midrule
\texttt{no\_\allowbreak cycle\_\allowbreak of\_\allowbreak
  window\_\allowbreak bound} &
  \texttt{[propext, Quot.sound]} \\
\texttt{failure\_\allowbreak window\_\allowbreak contradicts\_\allowbreak
  window\_\allowbreak corrected} &
  \texttt{[propext, Quot.sound]} \\
\texttt{five\_\allowbreak x\_\allowbreak plus\_\allowbreak one\_\allowbreak
  diverges\_\allowbreak at\_\allowbreak 7\_\allowbreak of\_\allowbreak
  window\_\allowbreak bound} &
  \texttt{[propext, Classical.choice, Quot.sound]} \\
\texttt{DwdbDiv.dwdbDivFinalAssembly} &
  \texttt{[propext, Classical.choice, Quot.sound]} \\
\bottomrule
\end{tabular}
\end{center}

\subsection{Acceptance criteria}

\begin{enumerate}[leftmargin=*]
\item No \texttt{sorry} in any theorem consumed by the main chain.
\item No custom \texttt{axiom}; only \texttt{propext},
  \texttt{Classical.choice}, and \texttt{Quot.sound} are allowed.
\item No \texttt{native\_decide} in the main chain; the finite base is
  expanded explicitly.
\item \texttt{lake build CycleBridge} and
  \texttt{lake build DwdbDiv} pass.
\item \texttt{unified\_\allowbreak core\_\allowbreak final\_\allowbreak
  no\_\allowbreak hge} is promoted only after its \texttt{sorry} is
  replaced by a proof.
\item The final theorem name and statement are exactly
  \begin{verbatim}
FinalStatement.five_x_plus_one_diverges_at_7 :
  IsUnboundedOrbit 7
  \end{verbatim}
\end{enumerate}

\subsection{Update workflow}

Whenever the Lean repository changes, the following steps are run:
\begin{enumerate}[leftmargin=*]
\item rebuild the affected targets;
\item re-run the axiom audit;
\item update the dependency tables in this chapter and in the main
  paper;
\item update the status ledger in \texttt{STATUS.md};
\item recompile the paper and the manual;
\item do not promote any pending statement before its ledger entry is
  changed.
\end{enumerate}
